% Original: PDF strana 25, gornji slajd; odobren prvi deo s049.
\slajdid{02-s049}
\begin{frame}[label=02-s049]{NILI realizacija: tri različite vrste čipova}
\setlength{\parskip}{0pt}
% element: 02-s049:e01
Za realizaciju NILI kolima po pravilu je najlakše poći od proizvoda zbirova:
\[\begin{aligned}
F&=\textcolor{DEcrvena}{(D+\bar B)(\bar D+\bar C)(\bar D+\bar A)}\\[1.5mm]
 &=\overline{\overline{(D+\bar B)(\bar D+\bar C)(\bar D+\bar A)}}
 =\overline{\overline{(D+\bar B)}+\overline{(\bar D+\bar C)}+\overline{(\bar D+\bar A)}}.
\end{aligned}\]
% element: 02-s049:e02
Potrebni su:
\begin{itemize}
\item jedan čip sa najmanje \textbf{četiri invertora};
\item jedan čip sa najmanje \textbf{tri dvoulazna NILI kola};
\item jedan čip sa najmanje \textbf{jednim troulaznim NILI kolom}.
\end{itemize}
Ukupno su to \textbf{tri čipa različitih vrsta}.
\par\vspace{2mm}
% element: 02-s049:e04
\Rezultat{Ni ova NILI transformacija \alert{\textbf{ne dokazuje minimalnost}} prema klasičnoj minimizaciji.}
\end{frame}
